% !TeX root = ../main.tex
% 原始章节：01-hardware-preparation.Rmd

\chapter{基础知识与技能准备}\label{chap:01-hardware-preparation}

\section{Linux环境与EDA工具掌握}\label{sec:01-hardware-preparation-001}

因为硬件开发的大多数工具仅能在Linux环境下运行，所以我们的大多数开发和测试工作也都在Linux环境下进行。我们认为有必要对一些Linux命令行界面下的基础技能进行简要介绍。本文选用的发行版Linux系统为Ubuntu22.04系统，后面的介绍也将基于该系统展开。

\subsection{Shell}\label{sec:01-hardware-preparation-002}

通常情况下，我们将操作系统的核心部分称为内核（英文：Kernel），而与之相对的最外层则被称为壳（英文：Shell）。Shell 是用于访问操作系统服务的用户界面。一般而言，操作系统的 Shell 可以使用命令行界面（CLI）或图形用户界面（GUI）。命令行界面是一种纯文本的界面，可以接收用户发送的命令，并在命令存在于环境中时提供相应的功能。

在 Ubuntu 操作系统中，我们默认使用的命令行界面 Shell 是 bash。bash 是一款基于 GNU 的免费开源软件。

打开 Linux 命令行界面后，您会首先看到光标前的以下内容：@ 符号前是用户名，@ 符号后是计算机名，冒号后显示当前所在的文件目录（其中 / 表示根目录，\textasciitilde{} 表示主目录，通常等同于 \texttt{/home/\textless{}user\_name\textgreater{}}）。最后的 \$ 或 \# 符号分别表示当前用户是普通用户还是超级用户 root。

% 代码语言：BookShell
\begin{CodeBlock}[language=BookShell]
chip-env@chipenv-virtual-machine:~$ 
chip-env@chipenv-virtual-machine:~$ sudo su
[sudo] password for chip-env: 
root@chipenv-virtual-machine:/home/chip-env# 
\end{CodeBlock}

面对全新的操作界面时，我们首先需要了解当前目录中包含哪些文件。这时，可以使用 \texttt{ls} 命令。该命令能够以彩色高亮显示文件，以区分不同类型的文件，是使用频率较高的命令。其输出信息的详细示例如下。

通常情况下，如果省略 \texttt{ls} 命令的参数，它会默认显示当前目录下的所有非隐藏文件。若要显示隐藏文件，需要添加 \texttt{-a} 选项；若要查看文件的详细信息，则需要添加 \texttt{-l} 选项。

% 代码语言：BookShell
\begin{CodeBlock}[language=BookShell]
ls - list directory contents
用法:ls [选项]... [文件]...
选项（常用）：
	-a 不隐藏任何以. 开始的项目
	-l 使用长格式显示文件详细信息

# 一般的Ubuntu系统中输入ls可看到如下文件夹
chip-env@chipenv-virtual-machine:~$ ls
Desktop  Documents  Downloads  Music  Pictures  Public  snap  Templates  Videos
\end{CodeBlock}

我们可以使用 \texttt{mkdir} 命令来创建文件目录（类似于 Windows 系统中的文件夹）。该命令的参数是新目录的名称，例如，\texttt{mkdir\ newdir} 将创建一个名为 \texttt{newdir} 的目录。

% 代码语言：BookShell
\begin{CodeBlock}[language=BookShell]
mkdir - make directories
用法:mkdir [选项]... 目录...
\end{CodeBlock}

删除空目录的命令是 \texttt{rmdir}。该命令的参数是需要删除的空目录的名称。请注意，只有空目录才能使用 \texttt{rmdir} 命令进行删除。\texttt{rmdir} 的基本用法如下：

% 代码语言：BookShell
\begin{CodeBlock}[language=BookShell]
rmdir - remove empty directories
用法:rmdir [选项]... 目录...
\end{CodeBlock}

如果目录非空，您需要使用 \texttt{rm} 命令。\texttt{rm} 命令不仅可以删除一个或多个文件，还可以删除目录及其下属的所有文件和子目录。对于符号链接文件，\texttt{rm} 仅删除链接文件本身，而保留原有文件不变。\texttt{rm} 命令的基本用法如下：

% 代码语言：BookShell
\begin{CodeBlock}[language=BookShell]
rm - remove files or directories
用法:rm [选项]... 文件...
选项（常用）：
	-r 递归删除目录及其内容
	-f 强制删除。忽略不存在的文件，不提示确认
\end{CodeBlock}

\begin{quote}
\textbf{Note:} 使用 rm 命令要格外小心，因为一旦删除了一个文件，就无法再恢复它。所以，在删除文件之前，最好再看一下文件的内容，确定是否真要删除。
\end{quote}

此外，\texttt{rm} 命令可以使用 \texttt{-i} 选项，这在删除多个文件时尤为有用。启用此选项后，系统会逐一确认是否删除每个文件。您需要输入 \texttt{y} 并按 Enter 键以确认删除。如果仅按 Enter 键或输入其他字符，文件将不会被删除。

相对的，\texttt{-f} 选项用于强制删除文件或目录，不会进行用户确认。结合 \texttt{-r} 选项使用，可以实现递归强制删除，这将删除指定目录及其所有文件和子目录，可能导致毁灭性效果。例如，\texttt{rm\ -rf\ /} 命令会强制递归删除整个文件系统中的所有文件。请务必小心，绝不要轻易尝试此操作！

\begin{quote}
\textbf{Note:} 在需要键入文件名/目录名时，可以使用 TAB 键补足全名，当有多种补足可能时双击 TAB 键可以显示所有可能选项。
\end{quote}

在学会创建目录后，我们需要进入新建的目录中以创建和修改文件。可以使用 \texttt{cd} 命令来切换到指定的工作目录 \texttt{dirname}，其中 \texttt{dirname} 可以是绝对路径或相对路径。如果省略目录名称，则会切换到用户的主目录（等同于执行 \texttt{cd} 命令）。

此外，在 Linux 文件系统中，\texttt{.} 表示当前目录，\texttt{..} 表示当前目录的上一层目录。在使用 \texttt{ls\ -a} 命令时，您会看到目录中包含 \texttt{.} 和 \texttt{..}，它们分别代表当前目录和上一级目录。

% 代码语言：BookShell
\begin{CodeBlock}[language=BookShell]
cd - change working directory
用法:cd [路径]
e.g.
	cd . 变更至当前目录
	cd .. 变更至父目录
\end{CodeBlock}

在进入新建的目录后，目录一般为空，我们可以通过以下方法创建一些文件：

\begin{enumerate}
\def\labelenumi{\arabic{enumi}.}
\tightlist
\item
  \textbf{重定向输出}：可以使用 \texttt{\textgreater{}} 操作符创建一个空文件。例如，\texttt{\textgreater{}\ filename} 将创建一个名为 \texttt{filename} 的空文件。
\item
  \textbf{文本编辑工具}：使用文本编辑工具如 \texttt{vim} 或 \texttt{nano} 来新建并保存文件。使用命令 \texttt{vim\ filename} 或 \texttt{nano\ filename} 可以直接进行编辑。由于这些工具允许编辑并保存文件，许多用户更倾向于使用它们。
\end{enumerate}

读者可能会对``重定向输出''感到好奇。重定向输出是指使用 \texttt{\textgreater{}} 操作符将前一个命令的输出数据重定向到指定的文件中。例如，命令 \texttt{ls\ /\ \textgreater{}\ filename} 会将根目录下的文件名输出到当前目录下的 \texttt{filename} 文件中。

此外，还有以下几种重定向操作： - \textbf{重定向输入 (\texttt{\textless{}})}：将指定文件中的数据作为输入传递给前一个命令。例如，\texttt{command\ \textless{}\ filename} 会将 \texttt{filename} 文件中的数据作为输入传递给 \texttt{command}。

另外，\texttt{echo} 命令用于在 Shell 中打印变量的值或指定的字符串。简单来说，\texttt{echo} 命令会将其参数输出到屏幕上。将 \texttt{echo} 命令与重定向相结合，可以实现将输出数据写入到文件中，具体效果可以自行尝试。

最后，要查看文件内容，可以使用 \texttt{cat} 命令。执行 \texttt{cat\ filename} 可以将文件内容输出到屏幕上。如果需要显示行号，可以使用 \texttt{-n} 选项，例如 \texttt{cat\ -n\ filename}。

% 代码语言：BookShell
\begin{CodeBlock}[language=BookShell]
cat - concatenate files and print
用法:cat [选项]... [文件]...
选项（常用）：
	-n 对输出的所有行编号
\end{CodeBlock}

之后就是对于文件的操作：复制和移动。复制命令为\texttt{cp}，命令的第一个参数为源文件路径，命令的第二个参数为目标文件路径。

% 代码语言：BookShell
\begin{CodeBlock}[language=BookShell]
cp - copy files and directories
用法:cp [选项]... 源文件... 目录
选项（常用）：
	-r 递归复制目录及其子目录内的所有内容
\end{CodeBlock}

移动文件的命令是 \texttt{mv}，其操作与 \texttt{cp} 命令类似。例如，\texttt{mv\ file\ ../file\_mv} 将当前目录中的 \texttt{file} 文件移动到上一层目录，并重命名为 \texttt{file\_mv}。

此外，您可能已经注意到，\texttt{mv} 命令也可以用于文件重命名。在 Linux 系统中，如果要重命名文件，只需使用 \texttt{mv\ oldname\ newname} 命令即可完成。

% 代码语言：BookShell
\begin{CodeBlock}[language=BookShell]
mv - move/rename file
用法:mv [选项]... 源文件... 目录
\end{CodeBlock}

在工作中，您可能会遇到需要多次使用单条或多条复杂命令的情况。初学者可能会选择将这些命令保存到一个文件中，然后复制粘贴执行。然而，其实可以通过批处理脚本来简化操作。

简单来说，批处理脚本是一个存储了单条或多条命令的文本文件。在 Linux 系统中，有一种快速执行批处理脚本的方法，这与 Windows 系统中的 \texttt{.bat} 批处理脚本类似------即使用 \texttt{source} 命令。\texttt{source} 命令是 bash 的内置命令，通常也可以用点命令 \texttt{.} 来替代。

这两个命令都接受一个脚本文件作为参数，并在当前 Shell 环境中执行该脚本，而不会启动新的子进程。脚本中设置的所有变量都将成为当前 Shell 环境的一部分。

% 代码语言：BookShell
\begin{CodeBlock}[language=BookShell]
source - execute commands in file
用法:source 文件名 [参数]
注：source 读取脚本文件，不要求文件具有执行权限
\end{CodeBlock}

此外，还有两种常用的查找命令：\texttt{find} 和 \texttt{grep}。

使用 \texttt{find} 命令结合 \texttt{-name} 选项，可以在当前目录及其子目录中递归查找符合指定文件名的文件，并将这些文件的路径输出到屏幕上。

% 代码语言：BookShell
\begin{CodeBlock}[language=BookShell]
find - search for files in a directory hierarchy
用法:find -name 文件名
\end{CodeBlock}

\texttt{grep} 是一种功能强大的文本搜索工具，能够利用正则表达式对文本进行搜索，并将匹配的行打印出来。简而言之，\texttt{grep} 命令可以从文件中查找包含指定模式（pattern）的行，并将匹配的行以及文件路径输出到屏幕上。当你需要在整个项目目录中查找特定的函数名、变量名等文本时，\texttt{grep} 是一个非常有效的工具。

% 代码语言：BookShell
\begin{CodeBlock}[language=BookShell]
grep - print lines matching a pattern
用法:grep [选项]... PATTERN [FILE]...
选项（常用）：
	-a 不忽略二进制数据进行搜索
	-i 忽略匹配时的大小写差异
	-r 从文件夹递归查找
	-n 显示行号
\end{CodeBlock}

以上内容介绍了 Linux 系统中的一些入门级常用操作命令及其常用选项。如果希望了解这些命令的其他功能选项，或获取新命令的详细说明，可以使用 Linux 系统下的帮助命令------\texttt{man} 命令。通过使用 \texttt{man} 指令，用户可以查看有关 Linux 指令、配置文件以及编程等方面的帮助信息。

% 代码语言：BookShell
\begin{CodeBlock}[language=BookShell]
man - manual
用法:man page
e.g.
	man ls
\end{CodeBlock}

最后，介绍几个常用的快捷键：

\begin{itemize}
\tightlist
\item
  \textbf{Ctrl+C}：终止当前程序的执行。
\item
  \textbf{Ctrl+Z}：挂起当前程序。
\item
  \textbf{Ctrl+D}：终止输入（在使用 Shell 时，则退出当前 Shell）。
\item
  \textbf{Ctrl+L}：清屏。
\end{itemize}

需要注意的是，使用 \textbf{Ctrl+Z} 挂起程序后，会显示该程序的挂起编号。如果要恢复该程序，可以使用 \texttt{fg\ {[}job\_spec{]}} 命令，其中 \texttt{job\_spec} 为挂起编号。如果不输入编号，则默认为最近挂起的进程。

\begin{quote}
\textbf{Note:} 在多数 shell 中，四个方向键也是有各自特定的功能的：← 和 → 可以控制光标的位置，↑ 和 ↓ 可以切换最近使用过的命令
\end{quote}

\subsection{Vim}\label{sec:01-hardware-preparation-003}

Vim 被誉为``编辑器之神''，是一款专为程序员设计的高效编辑器，特别适合代码编辑。对于习惯图形化界面的文本编辑软件的用户而言，初次接触 Vim 可能会感到不习惯和不顺手。为此，下面总结了一些常用的基本操作，以帮助大家快速入门。同时，推荐大家阅读一篇高质量的 Vim 教程------\href{http://coolshell.cn/articles/5426.html/}{《简明 Vim 练级攻略》}。只需十多分钟的阅读，你就可以掌握 Vim 的所有基本操作。当然，真正精通 Vim 需要较长时间的练习。如果你只打算将 Vim 作为记事本使用，那么几分钟的学习即可。其他相关资料网上有很多，可以随时查阅。

\subsection{Makefile}\label{sec:01-hardware-preparation-004}

当我们准备投入代码编写时，常常会面临一个难题：从哪里开始阅读这堆代码？答案是：从 Makefile 开始。在你感到无从下手时，通常从 Makefile 入手是一个不错的选择。那么，什么是 \texttt{make}？什么又是 Makefile 呢？

\texttt{make} 是一种工具，通常用于维护工程项目。它根据时间戳自动判断项目中哪些部分需要重新编译，只编译必要的部分。\texttt{make} 工具会读取 Makefile 文件，并根据 Makefile 中的内容执行相应的编译操作。Makefile 类似于之前接触过的 VC 工程文件，但不同于 VC 的图形界面，Makefile 通过类似脚本的方式实现功能。

与 VC 工程文件相比，Makefile 具有更高的灵活性（但这也意味着学习成本较高）。Makefile 可以方便地管理大型项目，并且理论上支持任意编程语言，只要编译器可以通过 shell 命令调用。当项目涉及多种语言且构建流程复杂时，Makefile 能够展示其强大的能力。

为了帮助你更清楚地理解 Makefile 的基本概念，我们将编写一个简单的 Makefile。假设我们有一个 Hello World 程序需要编译，我们将为它编写一个简易的 Makefile。首先，如果没有 Makefile，我们需要使用如下指令来编译这个程序：

% 代码语言：BookShell
\begin{CodeBlock}[language=BookShell]
# 直接使用 gcc 编译 Hello World 程序
$ gcc -o hello_world hello_world.c
\end{CodeBlock}

那么，如果我们想将编译过程写成 Makefile，该如何操作呢？Makefile 的基本格式如下所示：

% 代码语言：BookMake
\begin{CodeBlock}[language=BookMake]
target: dependencies
	command 1
	command 2
	...
	command n
\end{CodeBlock}

在 Makefile 中，\texttt{target} 是我们构建的目标，而 \texttt{dependencies} 是构建该目标所需的其他文件或目标。接着是生成目标所需执行的指令。需要特别注意的是，每一条指令（\texttt{command}）前必须有一个 TAB 键，而不能是空格，否则 \texttt{make} 将会报错。

我们的简易 Makefile 可以写成如下所示，执行 \texttt{make} 命令即可生成 \texttt{hello\_world} 这个可执行文件：

% 代码语言：BookMake
\begin{CodeBlock}[language=BookMake]
hello_world: hello_world.o
	gcc -o hello_world hello_world.o

hello_world.o: hello_world.c
	gcc -c hello_world.c

clean:
	rm -f hello_world hello_world.o
\end{CodeBlock}

在这个示例中：

\begin{itemize}
\tightlist
\item
  \texttt{hello\_world} 是最终生成的目标文件，依赖于 \texttt{hello\_world.o}。
\item
  \texttt{hello\_world.o} 是一个中间目标，依赖于源文件 \texttt{hello\_world.c}。
\item
  \texttt{clean} 是一个伪目标，用于清理生成的文件。
\end{itemize}

为帮助读者更好地学习 Makefile，这里提供两个新手友好的网址（当然我们默认读者有一定的英语阅读基础）：

\begin{enumerate}
\def\labelenumi{\arabic{enumi}.}
\item
  \href{http://www.cs.colby.edu/maxwell/courses/tutorials/maketutor}{Makefile 新手教程}：该网站详细描述了如何从一个简单的 Makefile 入手，逐步扩展其功能，最终编写出一个相对完善的 Makefile。建议同学们从这个网站开始，按照示例创建自己的 Makefile。
\item
  \href{http://www.gnu.org/software/make/manual/make.html\#Reading-Makefiles}{GNU Make Manual}：这个网站提供了全面的 Makefile 语法说明。如果在实验过程中遇到不熟悉的 Makefile 语法，可以在此网站上找到详细的解释和答案。
\end{enumerate}

\subsection{相关书籍和资料}\label{sec:01-hardware-preparation-005}

几本推荐入门阅读的Linux书籍或在线教程：

\begin{itemize}
\tightlist
\item
  《鸟哥的Linux私房菜-基础学习篇》：这本书主要介绍了Linux系统下的命令，并教授如何使用这个系统以及简单的管理系统，很适合初学Linux系统的人。
\item
  《Linux Shell脚本攻略》：这本书讲解了Shell脚本的概念和作用，它其实就是Linux命令的集合。虽然这本书的内容不是很简单，但在熟悉一些基本的命令和参数后，阅读起来会更加容易。
\item
  《Pratical Vim》：这本书是vim的入门经典之作，书中的内容以实用为主，通过简单的例子和解释来帮助读者快速掌握vim的使用方法。
\item
  《跟我一起写Makefile》：这本书娓娓道来，适合初学者，可以帮助你快速掌握Makefile的编写技巧。
\end{itemize}

\subsection{EDA工具掌握}\label{sec:01-hardware-preparation-006}

芯片设计EDA工具与流程学习可以参考我们提供的\href{https://www.bilibili.com/video/BV1gXe8eUEmd}{培训视频}，视频基于敏捷开发环境，介绍了包括基本的设计与开发流程、工具环境安装与使用基本操作、脚本环境搭建与流程测试等内容，可以一个简单的CPU核心或者模块来对工具环境进行验证，确保流程畅通。

\section{数字电路与逻辑设计}\label{sec:01-hardware-preparation-007}

设计过程中我们主要会用到Verilog和SystemVerilog语言进行开发，读者需要认真学习两种语言，这里推荐一些相关的学习教材或书籍：

\begin{itemize}
\item
  《Verilog HDL高级数字设计》：这本书详细介绍了如何使用Verilog HDL对数字系统进行建模、设计与验证，对ASIC/FPGA系统芯片工程设计开发的关键技术与流程进行了深入讲解，内容涵盖集成电路芯片系统的前后端工程设计与实现中的关键技术及设计案例。
\item
  《Verilog编程艺术》：也是比较经典的Verilog HDL书籍，更适合有一定基础的读者阅读，书中主要介绍了Verilog编程的方法论和实用性，深入探讨了编码风格、语言特性、简洁高效和时钟复位等实际问题，并深入探讨了如何避免使用易混淆和易错误的语句，如何避免前后仿真不一致，以及如何充分发挥Verilog-2001的特性等。
\item
  \href{https://hdlbits.01xz.net/wiki/Main_Page}{HDLBits}：这是一个在线练习题目网站，的练习题目涵盖了从基础到高级的各种主题，包括组合逻辑、时序逻辑、状态机设计等。通过完成这些练习，学习者可以提高他们的硬件描述语言编程能力，并且加深对数字电路设计的理解。
\end{itemize}

\section{硬件设计仿真环境搭建}\label{sec:01-hardware-preparation-008}

实验环境的相关配置在教程文档中已有介绍这里不过多赘述。
